% ── Formal Echo: Boole ───────────────────────────────────────
% Appears after Guyon's early convergence.
% Tone: compressed, exact, almost airless.

\begin{formalecho}{George Boole}{The Discrete Collapse}

He works in a room with very little in it.

This is not poverty. It is method. Every object not in use is a
hypothesis his attention must reject, and he is trying to eliminate
rejection from the process entirely. He wants operations that
succeed or fail absolutely and without remainder.

The algebra he is building has two values. This seems to him not a
limitation but a clarification: the world, examined correctly,
resolves into states that either obtain or do not. The difficulty
has always been the examination, never the resolution.

He writes $x(1-x)=0$ and looks at it for longer than such a short
expression should require. What it says is: a thing cannot be
simultaneously itself and not itself. What it implies is that any
system capable of holding a thing in both states has failed to
finish thinking.

He believes this. He has believed nothing more completely.

Outside his window---he has tried to stop looking at windows---the
world continues its habit of failing to resolve. Clouds that might
be rain. A dog that might be asleep. A sentence in a letter on his
desk that might be an accusation or might be a compliment,
depending on something he cannot fix in place long enough to
evaluate.

He turns back to the page.

The page at least will hold still. What he writes there is either
true or false and the system he is building will, eventually, say
which. He does not ask for more than that. He asks for exactly
that, in terms precise enough that nothing else can enter.

He writes: \textit{Let the universe of discourse be defined.}

He means: let us agree on what can be inside.

He does not ask what happens to what is outside. This is the
condition. This is also the cost.

\end{formalecho}
